\(\mathcal P_L\) is the set of one-unit laws \(P\) such that \(P_X(\mathcal D)=1\); for every \(t\in\{0,1\}\) and \(H\in(0,1/2]\), \(\mathbb P_P(T=t,\lVert X-b\rVert_2<H)\geq c_0H^2\); \((\mathcal X_{0,P},\mathcal X_{1,P})\) is a known Borel partition of \(\mathcal D\) and \(T=\mathbf1\{X\in\mathcal X_{1,P}\}\) almost surely; \(Y=TY(1)+(1-T)Y(0)\) almost surely; for each \(t\) and for \(P_X\)-almost every \(x\in\mathcal X_{t,P}\), the conditional law of \(Y(t)\) given \(X=x\) is supported on \([0,1]\) and has mean \(\mu_t(x)\); and the point-anchored envelope
\[
 |\mu_t(x)-\theta_t|\leq L\lVert x-b\rVert_2
\]
holds for every \(t\) and for \(P_X\)-almost every side-\(t\) point \(x\) of radius at most \(1/2\), where \(\theta_t\in[0,1]\) and \(\mu_t(b)=\theta_t\).  The sampling condition \(\mathrm{ass:iid\text{-}sampling}\) is imposed on the sample and is not a member property of a one-unit law.